\documentclass[11pt,a4paper,sans]{moderncv}

\moderncvstyle{casual}
\moderncvcolor{green}

\usepackage[utf8]{inputenc}
\usepackage[a4paper,top=1cm,bottom=2.5cm,left=2.1cm,right=2.1cm]{geometry}
\usepackage{enumitem}
\setlist[itemize]{topsep=0pt,partopsep=0pt,parsep=0pt,itemsep=1pt,leftmargin=1.5em}
\linespread{0.96}
\usepackage{xcolor}
\usepackage{eso-pic}

% moderncv casual renders a multi-line contact footer on every page; the
% default footskip is too small and triggers a fancyhdr warning. Give it the
% height it wants so contact info doesn't collide with the bottom margin.
\setlength{\footskip}{52pt}

% Source Code Pro as the monospace family — closest pdflatex-native stand-in
% for the homepage's JetBrains Mono. Used below for the name, section
% headings, and the contact block so the CV echoes the site's command-line
% accent.
\usepackage[scale=0.95]{sourcecodepro}

\renewcommand*{\namefont}{\ttfamily\Huge\mdseries\upshape}
\renewcommand*{\titlefont}{\ttfamily\normalsize\mdseries\upshape}
\renewcommand*{\sectionfont}{\ttfamily\Large\mdseries\upshape}
\renewcommand*{\subsectionfont}{\ttfamily\large\mdseries\upshape}
\renewcommand*{\addressfont}{\ttfamily\small\mdseries\upshape}
\renewcommand*{\hintfont}{\ttfamily\small}

% Solid emerald bullets, slightly smaller than the moderncv default (which
% uses the hollow \faCircle[regular] at \tiny).
\renewcommand*{\labelitemi}{\strut\textcolor{color1}{\raisebox{0.45ex}{\fontsize{5pt}{0}\selectfont\faCircle}}}

% Emerald + warm stone palette, aligned with the homepage dark theme.
% moderncv uses color0/color1/color2 internally; override the accent (color1)
% and keep body text on a warm stone tone.
\definecolor{color1}{HTML}{047857} % emerald-700
\definecolor{color2}{HTML}{57534E} % stone-600

% Thin emerald sidebar on the left edge of every page — echoes the homepage
% accent without touching the main text block.
\AddToShipoutPictureBG{%
  \AtPageLowerLeft{\color{color1}\rule{8pt}{\paperheight}}%
}

\name{Vojtěch}{Kubáč}
\title{Full-stack engineer · web, applied AI, C++}
\address{Pardubice}{Czech Republic}{}
\phone[mobile]{(+420)~737~059~153}
\email{vojta.kubac@gmail.com}
\homepage{vojtechkubac.eu}
\social[linkedin][www.linkedin.com/in/vojtech-kubac/]{Vojtěch Kubáč}
\social[github][github.com/VojtechKubac/]{VojtechKubac}
\photo[76pt][1pt]{photo}

\begin{document}
\makecvtitle
\vspace*{-8mm}

\section{Experience}
\cventry{2024--today}{Software engineer}{Ententee}{}{}{%
  Full-stack product delivery across healthcare engagements plus applied-AI
  research, alongside my role at CFD Support. Stack spans Python (Django, FastAPI),
  Java / Spring Boot, React / TypeScript, PostgreSQL, with scikit-learn and
  Unity / C\texttt{\#}.
  \begin{itemize}%
    \item \textbf{DrChrono ·} US medical app uplift: Django/GraphQL, React/TS.
    \item \textbf{ALF ·} medical-faculty app for attestations: Spring Boot, React/TS, PostgreSQL.
    \item \textbf{Ententee Hub ·} internal ops and hiring tooling: FastAPI, React/TS, PostgreSQL.
    \item \textbf{AI Lab ·} crop-analytics CV and robot-swarm sim PoCs; scikit-learn, Unity.
  \end{itemize}}
\cventry{2019--today}{Software engineer}{CFD Support}{}{}{%
  Develop and maintain multiple modules of TCAE, a fluid dynamics simulation application
  deployable on desktop machines or computational clusters. Built around
  ParaView/VTK and OpenFOAM. Core stack: C\texttt{++}23, Python, CMake, Qt, GitLab CI.
  \begin{itemize}%
    \item \textbf{TCFD ·} Computational Fluid Dynamics module: OpenFOAM, ParaView UI.
    \item \textbf{TFEA ·} FEA (solid mechanics) module: CalculiX, CFD coupling, ParaView UI.
    \item \textbf{TCAA ·} aeroacoustics: FW--H / Farassat 1A far-field, SPL/octaves.
    \item \textbf{TBASE ·} SQLite simulation DB: reproducible archive, ParaView IPC, surrogate modelling.
    \item \textbf{TMESH ·} meshing module from STL/STEP (NetGen, Gmsh).
    \item \textbf{TCAE DevOps ·} GitLab CI, CTest, headless Qt, Linux/Windows installers.
  \end{itemize}}
\cventry{2019--2020}{PhD researcher}{Charles University (MFF UK)}{}{}{%
  Research in the group of Christoph Allolio on continuum modelling of lipid membranes.
  \begin{itemize}%
    \item Ran molecular dynamics simulations (GROMACS) of lipid membranes.
    \item Worked on implementation of a continuum mechanics solver in C\texttt{++} for membrane remodelling.
  \end{itemize}}

\section{Education}
\cventry{2017--2020}{Master's}{Charles University}{Prague}{Mathematical Modeling in Physics and Technology}{%
  Thesis: \emph{Comparison of possible formulations of
  fluid-structure interactions with application in biomechanics},
  supervised by RNDr.~Jaroslav~Hron, Ph.D.}
\cventry{2018--2019}{Erasmus}{Technical University of Munich}{Germany}{Faculty of Mathematics}{Machine learning, deep learning, high-performance computing.}
\cventry{2014--2017}{Bachelor's}{Charles University}{Prague}{General Mathematics}{}

\section{Skills}
\cvitem{Programming}{C\texttt{++} (C\texttt{++}23), Python, TypeScript, Java, C\texttt{\#}, Bash}
\cvitem{Platform}{Qt, CMake, CTest, SQLite, GitLab CI, NSIS, Linux, Git}
\cvitem{AI / ML}{RAG / LLM integration; AI-assisted development (Claude Code, Cursor, Copilot); multispectral-UAV computer vision; surrogate modelling over design spaces}
\cvitem{Scientific}{OpenFOAM, CalculiX, NetGen, Gmsh, ParaView, VTK, FEniCS, GROMACS}
\cvitem{Web BE}{Django, Spring Boot, FastAPI, GraphQL, PostgreSQL}
\cvitem{Web FE}{React, Svelte, Tailwind CSS, Vite}
\cvitem{Simulation}{Unity, ORCA/RVO2, A* pathfinding, multi-agent coordination}
\cvitem{Spoken}{Czech (native), English (C1), German (B1)}

\end{document}
